% ---------------------------------------------------------------------
\section{Proofs}\label{app:proofs}
% ---------------------------------------------------------------------

This section proves the formal results of the main text and states the auxiliary lemmas that the main text cites.
All statements hold for every EV separately, and the index $i$ is dropped where no confusion arises.
Inequalities are exact in real arithmetic, and the implementation checks them with one tolerance of $10^{-9}$ in kW or kWh that the unused part of the margin $\varepsilon_q$ absorbs.

\subsection{Size independence of the attention policy}\label{app:equivariance}

\begin{lemma}[Equivariance and size independence]\label{lem:equivariance}
Let $\varpi$ permute the charger tokens and their masks while the global token stays fixed.
\begin{enumerate}
\item[(i)] The means $\mu_{\theta}$ and standard deviations $\sigma_{\theta}$ are permutation-equivariant, and the value $V_\omega$ is permutation-invariant.
\item[(ii)] Outputs of unmasked tokens and the value do not depend on the features or the number of masked tokens.
\item[(iii)] The number of parameters does not depend on $N$.
\end{enumerate}
\end{lemma}

\begin{proof}
Embeddings, layer normalisation, feed-forward blocks and the actor head act row-wise, so they commute with the permutation.
For one attention head, the key mask depends only on the key index, so permuting tokens and mask maps the mask matrix $M$ to $PMP^\top$ for the permutation matrix $P$ that fixes the global token.
Row-wise softmax satisfies $\mathrm{softmax}(PXP^\top)=P\,\mathrm{softmax}(X)P^\top$, so a masked attention head maps permuted inputs to permuted outputs, and composition gives (i) for the actor.
A masked key receives zero weight in every softmax, so its value never enters the output of another token, and together with the row-wise structure this property gives (ii).
Attention pooling sums over keys and is therefore invariant, and every weight matrix has a shape fixed by $(F_c,F_g,d_m,n_H,n_L)$ alone, so (iii) holds.
The safety layer is permutation-equivariant as well, because its problem is invariant under joint permutation of $p$, $x$, $l$ and $u$ and its minimiser is unique.
\end{proof}

The parameter counts are 610,370 for the actor and 676,481 for the critic, 1,286,851 in total for every $N$.
For comparison, the small ablation network ($d_m=64$, $n_L=2$, feed-forward width 256) has 228,067 parameters.
The actor of the MLP baseline with 200 slots has 792,976 parameters, a number that grows linearly with the number of slots.

\subsection{Reachable energy and thresholds}\label{app:reach}

\begin{lemma}[Reachable energy]\label{lem:reach}
Under A1, $\Psi_k$ is continuous and non-decreasing for every $k\ge0$.
Any admissible power sequence from $E(t)$ gives $E(t+k)\le\Psi_k(E(t))$, with equality for full-power charging.
An EV is individually feasible at step $t$, meaning that some admissible control meets \Eq{M-eq:requirement}, if and only if $E(t)\ge B(\tau(t))$.
\end{lemma}



\begin{proof}
A composition of continuous non-decreasing maps is continuous and non-decreasing, so $\Psi_k=\varphi^{k}$ inherits both properties from A1.
For the reach bound we use induction on $k$, where the case $k=0$ is trivial.
In the induction step, an admissible power $p\le\bar P(E)$ gives $E(t+1)=E+\Delta t\,g(p)\le E+\Delta t\,g(\bar P(E))=\varphi(E)$ because $g$ is increasing, and the induction hypothesis from $E(t+1)$ with the monotonicity of $\Psi_{k-1}$ gives
\[
E(t+k)\le\Psi_{k-1}\big(E(t+1)\big)\le\Psi_{k-1}\big(\varphi(E)\big)=\Psi_k(E).
\]
Full-power charging gives equality at every step.

For the threshold characterisation we show $\{E\in[E^{\min},E^{\sm{cap}}]:\Psi_k(E)\ge\Ereq\}=[B(k),E^{\sm{cap}}]$ by induction, where the case $k=0$ is the definition $B(0)=\Ereq$.
The set in \Eq{M-eq:thresholds} is non-empty because $\varphi(E^{\sm{cap}})=E^{\sm{cap}}\ge\Ereq\ge B(k-1)$.
In the induction step, $\Psi_k(E)=\Psi_{k-1}(\varphi(E))\ge\Ereq$ holds if and only if $\varphi(E)\ge B(k-1)$, and since $\varphi$ is continuous and non-decreasing, $\{E:\varphi(E)\ge B(k-1)\}$ is a closed up-set with minimum $B(k)$.
By the reach bound, some admissible control meets \Eq{M-eq:requirement} from $E(t)$ if and only if full-power charging does, that is, if and only if $\Psi_{\tau(t)}(E(t))\ge\Ereq$.
\end{proof}

For a certified table, the first two properties of \Eq{M-eq:certified} imply $\Bh\ge B$ componentwise, because $\varphi(\Bh(k))\ge\Bh(k-1)\ge B(k-1)$ gives $\Bh(k)\ge B(k)$ by the minimality in \Eq{M-eq:thresholds}.

\subsection{Projection (Proposition~\ref{M-prop:projection} of the main text)}\label{app:projection}

\begin{proof}
(i) If $p\in\mathcal K(t)$, then $l\le p\le u$ gives $l\le u$ and $\sum_il_i\le\mathbf 1^\top p\le\sum_iu_i$, and combined with $-\Crev\le\mathbf 1^\top p\le\Ctil$ these inequalities give the four conditions.
Conversely, $p(\upsilon)=l+\upsilon(u-l)$ stays in the box for $\upsilon\in[0,1]$, and $\mathbf 1^\top p(\upsilon)$ takes every value in $[\sum_il_i,\sum_iu_i]$, an interval that meets $[-\Crev,\Ctil]$ under the four conditions.

(ii) The problem in \Eq{M-eq:proj} minimises a strictly convex quadratic over a non-empty polyhedron, so its minimiser exists and is unique, and the KKT conditions are necessary and sufficient.
With multipliers $\mu^{+},\mu^{-}\ge0$ for the aggregate constraints and $\nu^{l},\nu^{u}\ge0$ for the box, they read
\[
\begin{gathered}
p-x+(\mu^{+}-\mu^{-})\mathbf 1-\nu^{l}+\nu^{u}=0,\qquad \mu^{+}(\mathbf 1^\top p-\Ctil)=0,\qquad \mu^{-}(\mathbf 1^\top p+\Crev)=0,\\
\nu^{l}_i(p_i-l_i)=0,\qquad \nu^{u}_i(u_i-p_i)=0.
\end{gathered}
\]
The function $\mathcal S$ is continuous, non-increasing and piecewise linear, with $\mathcal S(\lambda)=\sum_iu_i$ for $\lambda\le\min_i(x_i-u_i)$ and $\mathcal S(\lambda)=\sum_il_i$ for $\lambda\ge\max_i(x_i-l_i)$.
If $\mathcal S(0)>\Ctil\ge\sum_il_i$, the intermediate value theorem gives a root of $\mathcal S(\lambda)=\Ctil$, and every such root is positive; the case $\mathcal S(0)<-\Crev\le\sum_iu_i$ is symmetric with a negative root.
Now set $p=\clip(x-\lambda^{*},l,u)$, $\mu^{+}=\max(\lambda^{*},0)$ and $\mu^{-}=\max(-\lambda^{*},0)$, and choose the box multipliers coordinate by coordinate.
Free coordinates with $l_i<x_i-\lambda^{*}<u_i$ receive $\nu_i^{l}=\nu_i^{u}=0$, coordinates with $x_i-\lambda^{*}\le l_i$ receive $\nu_i^{l}=l_i-x_i+\lambda^{*}\ge0$ and $\nu_i^{u}=0$, and the remaining ones receive $\nu_i^{u}=x_i-\lambda^{*}-u_i\ge0$ and $\nu_i^{l}=0$.
Stationarity then holds coordinate-wise, and complementarity holds because $\lambda^{*}>0$ only when $\mathbf 1^\top p=\Ctil$ and $\lambda^{*}<0$ only when $\mathbf 1^\top p=-\Crev$.
For $\lambda^{*}=0$ the point with $\mathbf 1^\top p=\mathcal S(0)$ is feasible, so in all cases $p=p^{*}$.
If $\lambda_1<\lambda_2$ are both roots, every coordinate $\clip(x_i-\lambda,l_i,u_i)$ is non-increasing in $\lambda$ while the sums agree, so all coordinates agree.

(iii) If the aggregate constraint is inactive, locally $p_i=x_i$ for $i\in\mathcal F$ and $p_i\in\{l_i,u_i\}$ is constant otherwise, so the Jacobian is $D$.
If it is active with target $\vartheta\in\{\Ctil,-\Crev\}$, then $\sum_{i\in\mathcal F}(x_i-\lambda)+\sum_{i\notin\mathcal F}p_i=\vartheta$.
Hence $\partial\lambda/\partial x_j=\mathbb 1[j\in\mathcal F]/|\mathcal F|$ and $\partial p_i/\partial x_j=\mathbb 1[i\in\mathcal F]\big(\mathbb 1[i=j]-\mathbb 1[j\in\mathcal F]/|\mathcal F|\big)$.
If $\mathcal F=\emptyset$, the point $p$ is locally constant, and on the null set where $p^{*}$ is not differentiable, the matrices of the adjacent regions are elements of Clarke's generalised Jacobian.
\end{proof}

The Jacobian is symmetric and idempotent, and as a Euclidean projection onto a convex set, $\Pi_t$ is non-expansive.

\subsection{One-step invariance (Lemma~\ref{M-lem:onestep} of the main text)}\label{app:onestep}

\begin{proof}
The bound $\pmn\le l_i$ holds by \Eq{M-eq:bounds}, and for $l_i\le u_i$ note first that $\pmn\le0\le\bar P(E)$.
Moreover, $g^{-1}\big((\Bh(\tau-1)-E)/\Delta t\big)\le\bar P(E)$ is equivalent to $\Bh(\tau-1)\le\varphi(E)$, and A1 with \Eq{M-eq:certified} gives $\varphi(E)\ge\varphi(\Bh(\tau))\ge\Bh(\tau-1)$.
For the second claim, let $y=(\Bh(\tau-1)-E)/\Delta t$; since $l_i\ge g^{-1}(y)$ and $g$ is increasing, every $p\ge l_i$ satisfies $g(p)\ge y$, and A3 gives $E_i(t+1)\ge E+\Delta t\,g(p)\ge\Bh(\tau-1)$.
\end{proof}

When $\tau\ge2$, the EV is individually feasible at $t+1$ because $\tau_i(t+1)=\tau-1$, and when $\tau=1$, the lemma gives $E_i(d_i+1)\ge\Bh(0)\ge\Ereq_i$.

\subsection{Feasibility invariance (Theorem~\ref{M-thm:invariance} and Corollary~\ref{M-cor:local} of the main text)}\label{app:theorem}

\begin{proof}[Proof of Theorem~\ref{M-thm:invariance}]
We prove by induction over $t$ the invariant $\mathrm I(t)$: $E_i(t)\ge\Bh_i(\tau_i(t))$ for all $i\in\mathcal N(t)$.
An EV belongs to $\mathcal N(t)$ either because it arrives at $t$ or because it was in $\mathcal N(t-1)$ with $d_i\ge t$.
In the first case its part of the invariant is A4, in the second case the induction step applies, and since $\mathcal N(0)$ contains arrivals only, $\mathrm I(0)$ holds.

Assume $\mathrm I(t)$, so that every $i\in\mathcal N(t)$ is individually feasible and $l_i(t)$ is given by the first case of \Eq{M-eq:bounds}.
Lemma~\ref{M-lem:onestep} gives $\pmn_i\le l_i\le u_i$, and empty chargers have $l_i=u_i=0$.
By A5 and Proposition~\ref{M-prop:projection}(i), $\mathcal K(t)\neq\emptyset$, so $p(t)=\Pi_t(x(t))\in\mathcal K(t)$ for every finite $x(t)$, and claim (b) follows from $p_i(t)\in[l_i,u_i]\subseteq[\pmn_i(E_i(t)),\bar P_i(E_i(t))]$.
For claim (a), $\mathbf 1^\top p(t)\in[-\Crev(t),\Ctil(t)]$, and A3 gives
\[
\sum_i\tilde p_i(t)\in\Big[\sum_ip_i(t),\ \sum_ip_i(t)+n(t)\varepsilon_q\Big]\subseteq\big[-\Crev(t),\ \Ctil(t)+n(t)\varepsilon_q\big]=\big[-\Crev(t),C(t)\big].
\]
For each $i\in\mathcal N(t)$, Lemma~\ref{M-lem:onestep} gives $E_i(t+1)\ge\Bh_i(\tau_i(t)-1)$.
An EV with $d_i=t$ therefore departs with $E_i(d_i+1)\ge\Bh_i(0)\ge\Ereq_i$, so claim (c) holds for it.
Each EV with $d_i>t$ stays in $\mathcal N(t+1)$ with $\tau_i(t+1)=\tau_i(t)-1$, so its part of $\mathrm I(t+1)$ holds, and EVs that arrive at $t+1$ satisfy their part by A4.
All EVs depart before the episode ends, so claim (c) holds for every EV of the episode.
\end{proof}

\begin{proof}[Proof of Corollary~\ref{M-cor:local}]
The per-EV part of the induction uses only $p_i(t)\in[l_i(t),u_i(t)]$ during the stay of EV $i$, together with A4 and Lemma~\ref{M-lem:onestep}.
Fallback outputs satisfy $\pmn\le p\le u$, so $p_i\ge l_i$ implies $p_i\in[l_i,u_i]$, and on feasible steps the transformer part of the induction step applies unchanged.
In case (iii) of the fallback rule, Lemma~\ref{lem:fallback} gives $\mathbf 1^\top p=\Ctil\ge-\Crev$, and A3 transfers this sum bound to the realised powers.
\end{proof}

\textbf{A policy-independent sufficient condition for A5.}
The lower bound $l_i$ is non-increasing in $E$, because both terms of its first case are, and we let $l_i^{\max}(t)$ denote $l_i$ evaluated at $E=\Bh_i(\tau_i(t))$.
Suppose that, at every step, $\Crev(t)\ge0$, $\Ctil(t)\ge-\Crev(t)$ and $\sum_{i\in\mathcal N(t)}l_i^{\max}(t)\le\Ctil(t)$.
Under $\mathrm I(t)$, $l_i(t)\le l_i^{\max}(t)$ and $\sum_iu_i\ge0\ge-\Crev$, so A5 holds at $t$, and the induction step of Theorem~\ref{M-thm:invariance} yields $\mathrm I(t+1)$; A5 then holds for every policy.
However, $l_i^{\max}(t)=\bar P_i(\Bh_i(\tau_i))$ whenever $\Bh_i(\tau_i)>E_i^{\min}$, so the condition asks the transformer to serve all non-trivially constrained EVs at full power at the same time.

\subsection{Fallback rule and deficit bound}\label{app:fallback}

\begin{lemma}[Fallback]\label{lem:fallback}
In case (iii) of the fallback rule with $b\ge0$, the output satisfies $\pmn\le p\le l$ and $\mathbf 1^\top p=\Ctil$, so the transformer limit and the physical limits hold.
Every EV that is individually feasible at step $t$ and receives $p_i=l_i$ stays individually feasible.
In cases (ii) and (iv), the output minimises the transformer violation over the box $[\pmn,u]$, so no admissible command avoids that violation.
\end{lemma}

\begin{lemma}[Deficit bound]\label{lem:deficit}
Assume A1, A3 and A4, and assume that $\bar P_i$ is non-increasing in $E$.
In the EV2Gym instantiation, the total departure shortfall of the fleet in an episode is at most the deficit of \Eq{M-eq:deficit}.
\end{lemma}

Lemma~\ref{lem:taper} in \Sec{sec:ev2gym} shows that EV2Gym's taper satisfies the extra condition of Lemma~\ref{lem:deficit}.



\begin{proof}[Proof of Lemma~\ref{lem:fallback}]
In case (iii), $\sum_i(l_i-\pmn_i)>\Ctil-\sum_i\pmn_i=b$, so the budget is exhausted before every EV reaches $l_i$; hence $\mathbf 1^\top p=\Ctil$ and $p\le l$.
Outside case (i) we also have $\Ctil\ge-\Crev$, and an individually feasible EV with $p_i=l_i$ stays feasible by Lemma~\ref{M-lem:onestep}.
Over the box $[\pmn,u]$, the sum $\mathbf 1^\top p$ is smallest at $\pmn$ and largest at $u$, so the outputs of cases (iv) and (ii) minimise the commanded violation.
\end{proof}

\begin{proof}[Proof of Lemma~\ref{lem:deficit}]
Let $\tilde l_i(t)=l_i(t)$ if $E_i(t)\ge\Bh_i(\tau_i(t))$ and $\tilde l_i(t)=u_i(t)$ otherwise, and define the feasibility deficit $\Delta E^{\sm{def}}(t)=(\Bh(\tau(t))-E(t))^{+}$ with $\tau(d_i+1)=0$.
The shortfall $(\Ereq-E(d_i+1))^{+}$ is at most $\Delta E^{\sm{def}}(d_i+1)$ because $\Bh(0)\ge\Ereq$, and $\Delta E^{\sm{def}}(t_i^{\sm{arr}})=0$ holds by A4.
It therefore suffices to show $\Delta E^{\sm{def}}(t+1)\le\Delta E^{\sm{def}}(t)+\Delta t\,\big(g(\tilde l)-g(p)\big)^{+}$ at every step and to sum these increments over the stay.
If $E\ge\Bh(\tau)$, then $E+\Delta t\,g(l)\ge\Bh(\tau-1)$ by \Eq{M-eq:bounds}, and A3 gives $E(t+1)\ge\Bh(\tau-1)-\Delta t\,(g(l)-g(p))^{+}$.
When $E<\Bh(\tau)$ instead, $\tilde l=u$ and $E+\Delta t\,g(u)=\varphi(E)$, and \Eq{M-eq:certified} with the 1-Lipschitz property of $\varphi$ gives $\Bh(\tau-1)-\varphi(E)\le\varphi(\Bh(\tau))-\varphi(E)\le\Bh(\tau)-E=\Delta E^{\sm{def}}(t)$.
Hence $\Bh(\tau-1)-E(t+1)\le\Delta E^{\sm{def}}(t)+\Delta t\,(g(u)-g(p))$ in this case as well.
For the fleet bound with $g$ the identity, $p_i\ge\tilde l_i$ holds on feasible steps and $p=u$ holds in case (ii) of the fallback rule.
In case (iii), $p\le l$ and $\sum_i(l_i-p_i)=\sum_il_i-\Ctil$, and in case (iv), $p=\pmn$, so summing the per-step increments over all EVs gives \Eq{M-eq:deficit}.
\end{proof}

\subsection{Conformal departure guarantee (Proposition~\ref{M-prop:conformal} of the main text)}\label{app:conformal}

\begin{proof}
(i) By exchangeability, the rank of the test score $S_{n+1}=\hat q(\mathbf w_i)-d_i$ among the scores $S_1,\dots,S_{n+1}$ is uniform up to ties, so $\mathbb P(S_{n+1}\le\hat Q_{1-\alpha})\ge k/(n_{\sm{cal}}+1)\ge1-\alpha$ \cite{vovk2022algorithmic,lei2018distributionfree}.
The event $S_{n+1}\le\hat Q_{1-\alpha}$ is equivalent to $d_i\ge\hat q(\mathbf w_i)-\hat Q_{1-\alpha}$, and because $d_i$ is an integer, also to $d_i\ge\lceil\hat q(\mathbf w_i)-\hat Q_{1-\alpha}\rceil$.
Taking the maximum with $t_i^{\sm{arr}}$ changes nothing, because $d_i\ge t_i^{\sm{arr}}$ holds surely.

(ii) Consider the event that $d_i\ge\hat d_i^{\sm{lo}}$ and that $p_i(t)\ge l_i(t)$ at every step of the stay.
On this event, Lemma~\ref{M-lem:onestep} with $\tau^{\sm{lo}}$, applied up to $t=\hat d_i^{\sm{lo}}$, gives $E_i(\hat d_i^{\sm{lo}}+1)\ge\Bh(0)$.
For the remaining steps $\hat d_i^{\sm{lo}}<t\le d_i$ we have $\tau^{\sm{lo}}=1$, and the third property of \Eq{M-eq:certified} gives $\Bh(1)\le\max\{\Bh(0),E^{\min}\}\le E_i(t)$.
Lemma~\ref{M-lem:onestep} with $\tau=1$ then gives $E_i(t+1)\ge\Bh(0)\ge\Ereq$ again, so $E_i(d_i+1)\ge\Ereq$ holds on the event, and the union bound gives the stated probability.
If A5 holds almost surely, the second probability vanishes and the bound becomes $1-\alpha$.
Note that the bound does not hold conditionally on the event that A5 held, because that event may be correlated with $d_i$.
\end{proof}

\subsection{Upper bound of the perfect-information LP}\label{app:oracle}

\begin{proposition}[Upper bound]\label{prop:oracle}
\begin{enumerate}
\item[(a)] Let $\pi$ be any controller whose closed-loop trajectory satisfies \Eqs{M-eq:dynamics}{M-eq:limits}, satisfies \Eq{M-eq:transformer} at every step and has a total shortfall $\sum_i(\Ereq_i-E_i(d_i+1))^{+}\le z^{*}$. Then $J^{\pi}\le J^{\sm{LP}}$.
\item[(b)] In EV2Gym, let $J^{\sm{LP,ceil}}$ be the optimum of the same LP with $\Ereq_i$ and $E_i^{\min}$ lowered by $G_i^{\max}=0.01\,(d_i-t_i^{\sm{arr}}+1)$ kWh. Then $J^{\pi}\le J^{\sm{LP,ceil}}$ for every controller whose reported powers satisfy \Eq{M-eq:transformer} and whose realised total shortfall is at most $z^{*}$.
\end{enumerate}
\end{proposition}

\begin{proof}
(a) Set $p_i^{\pm}(t)=p_i(t)^{\pm}$, take $E_i(t)$ from the trajectory and set $z_i=(\Ereq_i-E_i(d_i+1))^{+}$.
With these values, the energy constraints of \Eq{M-eq:lp} hold by \Eq{M-eq:dynamics} and \Eq{M-eq:limits}, and the power bounds hold because $p_i\le\bar P_i(E_i)\le\Pnom_i$ and $p_i\ge\pmn_i\ge-\Pdis_i$.
In addition, the departure constraint holds by the definition of $z_i$, and the transformer constraint is \Eq{M-eq:transformer} itself.
Since $\sum_iz_i\le z^{*}$, the trajectory is feasible for stage 2, and its objective equals $J^{\pi}$.

(b) Let $\tilde p_i(t)$ be the reported powers and $E_i^{\sm{nom}}(t)=E_i^{\sm{arr}}+\Delta t\sum_{s<t}\tilde p_i(s)$ the energy that the billed powers alone would produce.
EV2Gym bills exactly $\tilde p$, so $J^{\pi}$ is the LP objective at $\tilde p$.
The rounding adds less than 0.01 kWh per active step, so $E_i^{\sm{real}}(t)-G_i^{\max}\le E_i^{\sm{nom}}(t)\le E_i^{\sm{real}}(t)\le E_i^{\sm{cap}}$ holds along the trajectory.
Both the lowered energy floor and the lowered departure constraint therefore hold for $E^{\sm{nom}}$ with $z_i=(\Ereq_i-E_i^{\sm{real}}(d_i+1))^{+}$, and the power bounds hold because $\tilde p_i\le\bar P_i(E^{\sm{real}}_i)\le\Pnom_i$.
From here on, the argument is as in (a).
\end{proof}

The main text uses $J^{\sm{LP,ceil}}$ as the reference of the optimality gap, because a feasible controller inside EV2Gym can exceed $J^{\sm{LP}}$ through the unbilled rounding gain.
A gap with respect to $J^{\sm{LP}}$ could then be negative, which would be meaningless for a bound.
\Tab{tab:ts1} in \Sec{sec:results-supp} reports the gap with respect to both optima.
